\section{Original arithmetic and the selected advance}
The predecessor constructs a nonzero top-socle product from the available
prime powers inside either cyclic subgroup of the dyadic unit group, then
uses one outside factor to reach the target \cite[Sections 2--4]{Prior}.
This continuation allows the outside factors to contribute directions as
well: the alternatives $q$ and $r$ give the inside ratio $q/r$, while the
factor $r$ remains in the original integer word. It is an actual exchange
inside the bounded exponent box, not an assertion that a negative prime
exponent is available. The complete two-prime box decomposes into the
retained exchange lines below. Their maximal capacities have closed formulas
at every exponent multiplicity.

The SplitZero convention remains
$G(A)=\{\tau\}\sqcup A^\bullet$, $e=0_A^\bullet\ne\tau$.
All original choices are retained with their integer coefficients. Reduction
modulo two is a specified observation: an even positive coefficient becomes
a supported zero, not an absent choice. We neither assign a socle after
calculating successful states nor infer their existence from an abstract
augmentation alone.

Fix
\[
 k\ge3,\quad m=2^k,\quad o=2^{k-2},\quad u=2^{2k-5},\qquad
 p\equiv1\pmod8,\quad p>2u,\quad N=p+4u=\prod_{q}q^{e_q}. \tag{A1}
\]
The factors $q$ are distinct actual primes; the $e_q$ are their actual
multiplicities. The structural results permit any integer $p$ in (A1).
Primality is checked separately in examples and the prime census.

The original middle states at this fixed seed are in bijection with
\[
 Q\mid N,\quad Q\equiv-1\pmod m,\qquad
 R=N/Q,\quad a=(p+R)/4.                                    \tag{A2}
\]
For completeness, $K(u)=\prod_{\ell^v\parallel u}\ell^{\lceil v/2\rceil}
=2^{k-2}$ and $u\mid a^2\iff K(u)\mid a$. Since $m\mid4u$, (A2) gives
$R\equiv-p\pmod m$ and hence that availability. Both $Q,R$ are $3\pmod4$,
so at least three; $N<3p$ implies both are less than $p$. Also
$R\mid p+4u$ implies $R\mid a+u$. Reversing these steps proves the inverse
from the original gate and availability conditions.

Retain the complete normalization
\[
 d=\gcd(a,u),\quad(h,r,s_0)=(d^2/u,u/d,a/d),\quad
 \lambda=(r+s_0)/R,\quad
 (x,y,z)=(a,phs_0\lambda,phr\lambda).                       \tag{A3}
\]
For $v_q(a)=E$, $v_q(u)=F\le2E$, the exponents of $h,r,s_0$ are
$E-|F-E|$, $(F-E)_+$, $(E-F)_+$. Thus they are integers,
$a=hrs_0$, $u=hr^2$, and $\gcd(r,s_0)=1$.
All are units modulo $R$, so $R\mid a+u=hr(r+s_0)$ makes $\lambda$ integral.
Multiplying the reciprocal sum by $phrs_0\lambda$ gives
$p\lambda+r+s_0=(p+R)\lambda=4hrs_0\lambda$.
The ordered inverse is
\[
 \boxed{u=\frac{pa^2}{Ry-pa}.}                              \tag{A4}
\]
These are inherited Type II coordinates, not a new parametrization.
The explicit crosswalk to Elsholtz--Tao \cite[Section 2, (2.13)--(2.22)]{ET}
is $(a_{ET},b_{ET},c_{ET},d_{ET},e_{ET},f_{ET})=(r,s_0,\lambda,h,R,Q)$;
their denominator map returns $(x,z,y)$. This last transposition is recorded,
not silently removed. No complete classification for composite $p$ is imported.

\section{Actual two-prime exchange lines and their full inverse}
Let $q,r$ be distinct primes, with exponent budgets $E,F$. Set
\[
 \ell_c=\max(0,c-F),\quad h_c=\min(E,c),\quad 0\le c\le E+F.
\]
The whole original box is the disjoint union of the anti-diagonals
\[
 (c,t)\longmapsto(i,j)=(\ell_c+t,c-\ell_c-t),\qquad
 0\le t\le h_c-\ell_c.                                    \tag{B1}
\]
The inverse is $c=i+j$, $t=i-\ell_c$. Therefore the exact integer identity is
\[
 \sum_{i=0}^{E}\sum_{j=0}^{F}[q]^i[r]^j
 =\sum_{c=0}^{E+F}[q]^{\ell_c}[r]^{c-\ell_c}
       \sum_{t=0}^{h_c-\ell_c}([q]/[r])^t.                 \tag{B2}
\]
The right side is first an identity of Laurent monomials: all completed
monomials have nonnegative exponents within the original budgets.
Its reduction in a unit group is a later map. Every original monomial occurs
once. Parallel lines with fixed $i-j=d$ have the analogous inverse
\[
 \ell_d=\max(0,-d),\quad h_d=\min(F,E-d),\quad
 (i,j)=(d+\ell_d+t,\ell_d+t),\quad 0\le t\le h_d-\ell_d.     \tag{B3}
\]
Their ratio is $qr$ instead of $q/r$.

Fix either
\[
 H_3=\langle3\rangle=\{x:x\equiv1,3\pmod8\},\qquad
 H_{-3}=\langle-3\rangle=\{x:x\equiv1,5\pmod8\},             \tag{B4}
\]
inside $(\mathbb Z/m)^\times$. The identity
$v_2(3^{2^j}-1)=j+2$ for $j\ge1$, by successive difference-of-squares
factorizations, proves that these have order $o$. Each contains $p$, not $-1$.
Write $H=H_b$ and $\chi_b=1$ on $H$, $-1$ on its other coset.
When $\chi_b(q)=\chi_b(r)$, both $q/r$ and $qr$ lie in $H$.
The constant monomial in (B1) or (B3) remains separately recorded.

\begin{theorem}[Optimal line capacities, including their parity]
Let $E,F\ge1$ and $v=\min(E,F)$. If both primes lie outside $H$, prescribe
parity $\epsilon\in\{0,1\}$ of the constant monomial. The largest possible
parameter lengths in (B1) and (B3), respectively, are
\[
 \boxed{
 L_{\rm exchange}=v-\mathbf1_{E=F,\ v\not\equiv\epsilon\ (2)},\qquad
 L_{\rm parallel}=v-\mathbf1_{E=F,\ \epsilon=1}.}            \tag{B5}
\]
If both primes lie inside $H$, both maxima are $v$, with constant in $H$.
\end{theorem}
\begin{proof}
The anti-diagonal length is $\min(c,E,F,E+F-c)$, whose maximum $v$ occurs
precisely for $v\le c\le\max(E,F)$. If the budgets differ, this interval
contains both parities. If they agree, its only point is $c=v$, and either
neighbor gives length $v-1$. The parity of its constant is $c\pmod2$.
For parallel lines, difference zero gives length $v$ with even constant.
When the budgets differ, difference $1$ or $-1$ toward the larger budget also
gives length $v$ with odd constant. With equal budgets any odd difference
loses at least one, and difference one attains that loss. The same proof
without the parity constraint applies to inside primes.
\end{proof}

An explicit maximizing anti-diagonal uses $c=v$ when its parity is wanted,
$c=v+1$ when the budgets differ and its parity is not wanted, and $c=v-1$
otherwise. The parallel odd line uses offsets $(1,0)$ when $E>F$, and
$(0,1)$ otherwise. The formulas include length zero; no binomial at an
identity ratio is selected later.

\section{Affine packets from disjoint actual prime blocks}
Partition the actual prime indices into singletons and pairs. Pairs must
have the same $\chi_b$ value. For each block $B_i$ choose a retained affine
line
\[
 v_q(W)=a_q+\sigma_qt_i\quad(q\in B_i),\qquad 0\le t_i\le L_i,   \tag{C1}
\]
where the integer endpoint inequalities put every exponent in $[0,e_q]$.
The block step vector $(\sigma_q)_{q\in B_i}$ is nonzero. Use (B1) or (B3) for pairs. For a singleton
inside $H$, use exponent $t$. For a singleton outside $H$, choose exponent
$\epsilon+2t$, $0\le t\le\lfloor(e_q-\epsilon)/2\rfloor$.

Keep
\[
 A=\prod_q q^{a_q},\qquad
 g_i=\prod_{q\in B_i}q^{\sigma_q}\in H\pmod m,
 \qquad \chi_b(A)=-1.                                    \tag{C2}
\]
The $g_i$ may be positive rationals, but
$W(\mathbf t)=A\prod_i g_i^{t_i}$ is always an actual positive integer divisor
of $N$. For a fixed set of block lines this map is injective: in each block,
a nonzero step coordinate gives
\[
 t_i=(v_q(W)-a_q)/\sigma_q,
\]
and all other block coordinates must agree. Distinct blocks have disjoint
prime indices. The full inverse image is thus specified, not just its residue.
The omitted part of the original exponent box is retained as the complement
of this explicitly defined subset. Reusing one prime in two blocks would
invalidate this inverse and is not allowed.

Write $g_i=b^{d_i}\pmod m$, with $0\le d_i<o$. Discard identity steps only
from the selected packet by setting their selected count to zero; their
constant monomials remain in $A$. For nonidentity steps set
$w_i=2^{v_2(d_i)}$ with capacity $L_i$.
Optionally include one formal phase step $p^{-1}=b^{-s}$, where $p=b^s$;
if $s\ne0$ it has capacity one and weight $2^{v_2(s)}$.
It records residual/cofactor orientation and is not a factor of $N$.

\begin{lemma}[Inherited bounded selection and its inverse]
For powers-of-two weights below $o=2^n$ and capacities $L_i$, a choice
$0\le c_i\le L_i$ with $\sum_i c_iw_i=o-1$ exists exactly when
\[
 \boxed{\sum_{w_i\le2^j}L_iw_i\ge2^{j+1}-1
 \quad(0\le j<n).}                                       \tag{C3}
\]
A highest-weight-first greedy choice constructs it.
\end{lemma}
\begin{proof}
Reducing the total modulo $2^{j+1}$ proves necessity. For sufficiency,
select one unit, pair the remaining units into twos, and divide all remaining
weights by two. The new prefix mass is
$\lfloor(C_{j+1}-1)/2\rfloor\ge2^{j+1}-1$; induction proves existence.
All carries retain the original labels. To justify the greedy inverse,
among smaller dyadic coins total mass at least $w$ contains exactly $w$:
pair upwards, or else the remaining unpaired coins have total at most $w-1$.
Exchange such a submultiset for another available weight-$w$ coin until the
greedy amount is attained. Continue downwards.
\end{proof}
This lemma and its top-socle consequence are inherited from the pinned
predecessor \cite[Section 3, Lemma 1 and Theorem 2]{Prior}; their short proofs
are retained to make the present application standalone.

\begin{theorem}[Actual factor-exchange forcing]
Assume (A1), (C1)--(C2), and the explicit prefixes (C3), including the optional
phase. Then an original positive integral middle witness exists. All selected
integer exponents and both possible roles have terminating inverses.
Every predecessor single-chart socle certificate is contained in this class.
\end{theorem}
\begin{proof}
In $\mathbb F_2[C_o]=\mathbb F_2[T]/T^o$, $T=X+1$, an actual nonzero log
$d=2^vr$, $r$ odd, gives
$1+X^d=T^{2^v}U_d(T)$ with $U_d(0)=1$.
Choose $c_i$ of total weight $o-1$. Then
\[
 \prod_i(1+X^{d_i})^{c_i}=T^{o-1}=\sum_{j=0}^{o-1}X^j.     \tag{C4}
\]
The extra unit disappears above the top degree. Repeated squaring proves
the last equality.
Let $t\preceq c$ mean binary submask containment. Frobenius gives
$(1+Y)^c=\sum_{t\preceq c}Y^t$ in characteristic two. The corresponding
integer packet is therefore
\[
 \mathcal P(X)= (1+X^{-s})^\delta
      \prod_{i}\left(\sum_{t_i\preceq c_i}X^{d_it_i}\right),
 \qquad \delta\in\{0,1\}.                                \tag{C5}
\]
Every coordinate is an actual allowed $t_i\le c_i\le L_i$, once, rather
than a coloured subset of indistinguishable prime copies. By (C4), every
integer cyclic coefficient is odd and positive. Choose its target
$-A^{-1}\in H$. It returns
\[
 W(\mathbf t)p^{-\varepsilon}\equiv-1\pmod m,
 \qquad \varepsilon\in\{0,\delta\}.                        \tag{C6}
\]
For $\varepsilon=0$, take $Q=W$; for $\varepsilon=1$, take $R=W$ and
$Q=N/W$. These satisfy (A2), and (A3)--(A4) finish the return.

To find a word, factor (C5) into its binary binomials and retain their suffix
products modulo two. At an odd target, exactly one of the two next suffix
coefficients is odd. Follow it and subtract its exponent when selected.
The terminal target is zero. An even coefficient on the other branch is
not called absent. Given a returned state and its role, read $W=Q$ or $W=R$,
then invert the original affine lines and check the submask conditions.
For a fixed packet and role this is injective; after forgetting the role
there are at most two inverse words, $Q$ and $N/Q$.

For predecessor inclusion choose all singletons, exponent zero offsets inside,
one outside prime with offset one, and all other outside primes with offset
zero. Its inside steps and the phase are exactly the old items; set the extra
square-step counts to zero. The old selected vector proves (C3), and the base
is the old actual outside prime.
\end{proof}

The perfect-cover framework is classical \cite[Section 1; Theorem 12]{GH}.
A packet here may have different odd integer multiplicities at different
residues; it is not asserted to be a perfect uniform cover. The reduction and
the original coefficient vector are both retained. In particular (C4) has
augmentation zero but is nonzero, and every original cyclic coordinate is
occupied. This is a supported-zero observation, not a target-absence test.
A ratio $q/r$ is not declared a freely available arithmetic factor: its
compensating affine constant in (C1) travels through every calculation.

\paragraph{Original count, disjoint free factors, and collisions.}
Let $M_0$ consist only of primes on which all selected physical words have
exponent identically zero. Partially unused budgets on other primes are
not put into this separate counting factor. Put
\[
 A_+(M_0)=\#\{t\mid M_0:\chi_b(t)=1\}
 =\tfrac12\left[\prod_{q^e\parallel M_0}(e+1)+
 \prod_{q^e\parallel M_0}\sum_{j=0}^{e}\chi_b(q)^j\right].   \tag{C7}
\]
For each such $t$, use target $-(At)^{-1}$ in (C5).
Unique factorization separates $t$ from the physical packet word within
one role. Each original cofactor has at most two roles. Hence
\[
 \boxed{C_k(p)\ge A_+(M_0)\ (\delta=0),\qquad
 C_k(p)\ge\lceil A_+(M_0)/2\rceil\ (\delta=1).}             \tag{C8}
\]
All collisions on forgetting the role are exactly pairs of words whose
product is $N$. Neither duplicated roles nor equal residues are counted as
new original states.

\paragraph{All multiplicities and the declared search class.}
For each prime matching and each coset parity/mode choice, (B5) gives the
largest capacity of that type. Replacing another line by this maximal one
preserves its cyclic step and constant coset, and only increases the prefixes.
Since (C5) covers the entire needed coset, changing the numerical constant
within that same coset introduces no new target restriction. Thus the
canonical search is complete within the declared singleton and disjoint
parallel/exchange-line class. It is not complete for all ES states or all
possible subfamilies of the exponent box.
With $r$ distinct primes there are at most $r!$ such matchings and at most
$2^r$ mode/parity choices per matching. There are two charts.
The decision search consequently uses at most $2^{r+1}r!$ prefix tests,
each on $O(r+k)$ data. Factoring $N$ and expanding $o$-term suffix products
have their separate costs. These are finite bounds, not a fast-factorization
claim. The formulas for maximal lines do not enumerate huge prime powers.

\section{A sharp close-pair theorem and exact counts}
A particularly useful exchange step is the half-turn. Let $k\ge4$ and
$L=o/2$. Let $b$ be an actual prime congruent to $3$ or $5\pmod8$, and write
$H=\langle b\rangle$, which is one of (B4). Indeed $v_2(b^2-1)=3$,
and successive difference-of-squares factorizations give
$v_2(b^{2^j}-1)=j+2$ for $j\ge1$, proving its full order. Retain distinct outside primes
$q,r$ with
\[
 \boxed{q/r\equiv b^L=1+2^{k-1}\pmod{2^k}.}                \tag{D1}
\]
Write $N=b^e q^{a_q}r^{a_r}M$, where $a_q,a_r\ge1$,
$\gcd(bqr,M)=1$ and $e\ge0$. The value $e=0$ is allowed for the obstruction
statement; the sufficient threshold below is positive.

\begin{theorem}[Half-turn exchange, sharp uniform capacity]
The original subfamily
\[
 Q=b^i t q\quad\text{or}\quad Q=b^i t r,
 \qquad 0\le i\le e,\quad t\mid M,\quad t\in H
\]
has count $C_{\rm ex}$ satisfying
\[
 \boxed{\left\lfloor\frac{e+1}{L}\right\rfloor A_+(M)
 \le C_{\rm ex}\le
 \left\lceil\frac{e+1}{L}\right\rceil A_+(M).}              \tag{D2}
\]
If $L\mid e+1$, both bounds are equal. In particular $e\ge L-1$ forces an
original witness, with no further phase assumption on $p$ beyond (A1).
If $a_q=a_r=1$ and every prime factor of $M$ is inside $H$, this subfamily
is the entire original fixed-seed branch.
The uniform sufficient exponent $L-1$ is optimal over the stated integer
parameters.
\end{theorem}
\begin{proof}
For each retained $t$, define $0\le j_t<o$ by
$b^{j_t}=-(qt)^{-1}\pmod m$. The $q$ words hit precisely when
$i\equiv j_t\pmod o$; by (D1), the $r$ words hit precisely when
$i\equiv j_t-L\pmod o$. Together this is
$i\equiv j_t\pmod L$, with exactly one partner for each such $i$.
Thus the exact per-divisor count is
\[
 \max\left\{0,1+\left\lfloor\frac{e-(j_t\bmod L)}L\right\rfloor\right\}.
                                                               \tag{D3}
\]
This is between the two integer bounds in (D2). The maps are injective because
$q,r,b$ and the prime factors of $t$ are distinct. If the outside primes occur
exactly once and $M$ has no outside prime, any target cofactor must contain
exactly one of $q,r$, proving the completeness assertion.

For sharpness fix $b=3$, choose primes
$q\equiv-3$, $r\equiv-3^{L+1}\pmod m$, and put
$e=L-2$, $N=3^e qr$, $p=N-4u$.
The two prime classes are distinct and reduced. Dirichlet's theorem gives
arbitrarily large actual choices \cite[Theorem 18.1]{Sutherland}; $p$ is
an integer, not asserted prime in this general construction.
It satisfies $p\equiv1\pmod m$ and eventually (A1).
Every outside divisor is $3^iq$ or $3^ir$ with $0\le i\le L-2$.
The target requires $i\equiv L-1\pmod L$, so the complete branch is empty.
\end{proof}

This inverse can be computed without expanding $3^e$: find $j_t$ by binary
lifting in the cyclic group, reduce it modulo $L$, choose the appropriate
partner, and count the remaining occurrences using (D3). The count is exact
for arbitrarily large binary-encoded $e$. If $q/r\equiv1\pmod8$ but is not
the identity modulo $2^k$, its exchange weight in either chart is
\[
 \boxed{2^{v_2(q-r)-2}.}                                   \tag{D4}
\]
Indeed $v_2(q/r-1)=v_2(q-r)$ because $r$ is odd, and the nonzero even-log
identity from (B4) gives $v_2(b^j-1)=v_2(j)+2$.
Thus an actual difference between two available primes supplies a high
nilpotent layer. It is not an assumed new prime divisor.

\section{Exact strict examples and the required negative example}
At the certified hard prime
\[
 p=837601\equiv121\pmod{840},\quad k=5,\quad u=32,\quad
 N=837729=3^3\cdot19\cdot23\cdot71,                         \tag{E1}
\]
the two outside factors satisfy $23/71\equiv17=3^4\pmod{32}$.
Retain the exchange $71(23/71)^t$, $t=0,1$, and the three available factors
of $3$. The cyclic packet is
\[
 (1+X)(1+X^2)(1+X^4)=\sum_{j=0}^7X^j.                     \tag{E2}
\]
The integer constant is $71$, not one; its omitted use would invalidate
the original word. The target logarithm is $2$, giving $Q=3^2\cdot71=639$
and $R=1311$. The original marking is
\[
 (a,h,r,s_0,\lambda)=(209728,32,1,6554,5),
\]
\[
 \boxed{\frac4{837601}=\frac1{209728}+
       \frac1{878341912640}+\frac1{134016160}.}              \tag{E3}
\]
The unused prime $19$ is inside $H_3$, so (C8) gives at least two original
states. Here (D2) is an equality: the complete branch has exactly
$Q=639,1311$. The other ordered triple is
$(209560,1799158571990,274733128)$.
The two prior chart-prefix tests and the $3,11$ capacity test fail;
neither $Q$ nor $R$ has at most two prime occurrences. This is a strict
improvement against their actual union, not only one selected baseline.

At $p=1740481$, $k=5$,
$N=3^5\cdot13\cdot19\cdot29$ and $13/29\equiv17\pmod{32}$.
The same exchange proves positivity. Its cofactors are exactly $351,4959$;
one return is
\[
 (u,R,a,h,r,s_0,\lambda)=(32,4959,436360,2,4,54545,11),
\]
with ordered denominators $(436360,2088559795190,153162328)$.

The third strict added branch uses a different original line:
\[
 p=1794769,\quad k=5,\quad N=3^2\cdot13\cdot23^2\cdot29.
\]
Use the parallel line $(v_{23},v_{29})=(1+t,t)$, $0\le t\le1$.
Its constant is $23$ and its step $23\cdot29\equiv3^3\pmod{32}$.
Together with selected count two for the $3$ block and the phase $p^{-1}$,
the weights are $2,1,4$, of total seven. The inverse selects physical word
$207=3^2\cdot23$ in the residual role, giving the unique cofactor
$8671=13\cdot23\cdot29$. Both occurrences of $23$ were available; neither
was used twice by two nominally separate blocks.

The threshold really matters. At the certified hard prime
\[
 p=278881\equiv1\pmod{840},\quad k=5,\quad
 N=3^2\cdot29\cdot1069,                                   \tag{E4}
\]
we have $29/1069\equiv17\pmod{32}$, but $e=2=L-2$.
The complete fixed-seed branch is empty, exactly as in the sharp construction.
This is not an ES counterexample: at seed $u=2$, $Q=39$, $R=7151$ the
ordered denominators are $(71508,99711112740,2788810)$.

The new class does not subsume the earlier two-block theorem. At the inherited
$p=482187176641$, $k=6$, $N=3^7 11^5 37^2$, that theorem applies while no
certificate in the declared affine-line search class is found. Exhaustion
here is of the finite maximal-line class, not of all possible algebraic packets.

\section{Finite comparison, preserved limitations, and next arithmetic step}
The complete inherited test universe is the 4519 primes through two million
in the six residue classes modulo $840$:
\[
 1,121,169,289,361,529.
\]
It uses every valid $k\ge4$ with $p>2u_k$. It has 37289 branches and 4583 original middle states.
\begin{center}
\begin{tabular}{lrr}
\toprule
Criterion&Branches&Primes\\\midrule
Previous cyclic socle&1139&1119\\
Affine-line socle here&1443&1362\\
Prime/two-prime words in both roles&3051&2387\\
Short words, prior $3,11$ test, or previous socle&3068&2399\\
That union or the present criterion&3071&2401\\
All original divisors in this seed family&3097&2411\\\bottomrule
\end{tabular}
\end{center}
The three added branches are $(837601,5)$, $(1740481,5)$ and $(1794769,5)$.
The additional primes relative to the full previous union are precisely
$837601,1740481$. These are improvements to sufficient tests, not new overall
ES verification coverage. A census of positive rows is not a universal theorem.

The verifier checks all original ordered states in the census, and reconstructs
an actual word and count bound for every new certificate. Separate structural
checks include 81 complete two-prime boxes, 1364 dyadic capacity profiles,
526 selected cyclic packets, and all half-turn intervals at cyclic orders
$4,8,16,32,64,128$ in the specified ranges. A separate implementation uses
literal diagonal enumeration, direct residue tables and bounded knapsack.
It imports none of the main verifier. The eight false-transform controls
include deleting the exchange constant, reusing a prime in two blocks, and
replacing the signed exchange slope without moving its logarithm.

A failed arithmetic branch must now avoid every available affine-block prefix
certificate, in addition to the previous actual-stabilizer constraints. For
example, any available half-turn outside pair with a full-order prime block
forces the corresponding exponent below $L-1$ if that branch is empty.
The pair difference in (D4) and the actual exponent inventory are both
quantitative inputs. We have not shown that the related integers $p+4u_k$
for one prescribed prime cannot simultaneously avoid these conditions.
The missing assertion is about their factorization correlations, not about
whether the marked return exists after a certificate has been obtained.

A scoped primary-source search did not establish an earlier identical combined
application. That is not an exhaustive historical-priority determination.
The cyclic algebra and perfect-cover background are classical; the original
ES normalization and earlier top-socle lemma retain their attribution.
No new Lean build, independent mathematical review, analytic theta-norm
identification or global ES theorem is asserted.
